\begin{abstract}
    Spatial reasoning, the ability to determine where objects are, how they relate, and how they move in 3D, remains a fundamental challenge for vision-language models (VLMs).
    Tool-augmented agents attempt to address this by augmenting VLMs with specialist perception modules, yet their effectiveness is bounded by the \emph{action interface} through which those tools are invoked.
    In this work, we study how the design of this interface shapes the agent's capacity for open-ended spatial reasoning.
    Existing spatial agents either employ single-pass code execution, which commits to a full analysis strategy before any intermediate result is observed, or rely on a structured tool-call interface that often offers less flexibility for freely composing operations or tailoring the analysis to each task. Both designs offer limited flexibility for open-ended, complex 3D/4D spatial reasoning.
    We therefore propose \textbf{\ours{}}, a training-free framework for spatial reasoning that adopts code as the action interface.
    \ours{} maintains a stateful Python kernel pre-loaded with input frames and a suite of perception and geometry primitives, letting a VLM-backed agent write one executable cell per step conditioned on all prior outputs, enabling the agent to flexibly compose and manipulate perception results and adapt its analysis to both intermediate text and visual observations and the demands of each problem.
    Evaluated across \textbf{20 spatial reasoning benchmarks} spanning a broad range of static and dynamic 3D/4D spatial reasoning tasks, \ours{} achieves 59.9\% average accuracy, outperforming the recent spatial agent by \textbf{+11.2 points}, with consistent gains across six VLM backbones from two model families \emph{without any benchmark- or model-specific adaptation}.
\end{abstract}
